\section{Текущее подтверждённое состояние}

На полном проходе по всем 560 target linear тензорам модели получено:

\begin{itemize}[leftmargin=1.5em]
\item mean relMSE $= 3.91 \times 10^{-6}$;
\item median relMSE $= 2.81 \times 10^{-6}$;
\item p90 relMSE $= 3.07 \times 10^{-6}$;
\item max relMSE $= 4.26 \times 10^{-5}$.
\end{itemize}

На полном raw WikiText-2 gate:

\begin{center}
\begin{tabular}{lrr}
\toprule
Модель & PPL & Комментарий \\
\midrule
dense & 3.4304 & базовый воспроизводимый прогон \\
routed & 3.4350 & gap $= 0.0046$ \\
\bottomrule
\end{tabular}
\end{center}

Этот результат означает, что quality-side для текстовой метрики практически
закрыт: разрыв к dense очень мал.

Для дополнительной sanity-проверки был также выполнен более короткий raw
WikiText-2 gate на первых $100$ окнах. Он дал dense PPL $= 3.2749$ и routed
PPL $= 3.2773$, то есть gap всего $0.00238$. На этом же прогоне были измерены
простые eval-throughput числа: `52736` токенов обрабатывались со скоростью
примерно `1006.4` tok/s для dense и `1016.1` tok/s для routed. Это не decode
benchmark, но как короткая full-model sanity check эта метрика полезна: на
таком промежуточном gate routed path не показывает явного regress по качеству
или по общему evaluation throughput.

На полном HumanEval gate ($164$ задачи):

\begin{center}
\begin{tabular}{lrr}
\toprule
Модель & pass@1 & elapsed \\
\midrule
dense & 0.7317 & 1337.0 s \\
routed & 0.7195 & 1169.0 s \\
\bottomrule
\end{tabular}
\end{center}

Абсолютный gap по pass@1 составил $-0.0122$ в пользу dense-базовой модели.

По памяти на этом gate:

\begin{itemize}[leftmargin=1.5em]
\item eval peak allocated: примерно 44.6 / 45.9 / 44.6 GB на трёх видимых GPU;
\item route surgery peak allocated: примерно 70.2 / 71.5 / 70.2 GB.
\end{itemize}